\subsection{Notación del límite y del tablero}\label{sec:notacion}
En la tabla~\ref{tab:notacion} se distinguen los objetos del tablero finito de las cantidades
del perfil. En particular, un par marcado y la matriz de la que se toma su menor aparecen
juntos en el apéndice~\ref{app:boundary}.
\begin{center}
\small
\begin{tabular}{@{}L{.21\linewidth}L{.72\linewidth}@{}}
\toprule Símbolo & Significado\\ \midrule
\(\Phi,\ C_{\rm pred}\) & Coeficiente de \(n\) y término constante del predictor original.\\
\(R_Q,\ \cerr\) & Diferencia de los dos logaritmos y límite de \(nR_Q(n)\).\\
\(m,\ j,\ x\) & Filas vacías en un paso, tamaño de llegada \(\lfloor bn\rfloor\) y
 fracción residual \(m/n\). En los lemas de torres \(k\) es el tamaño de una máscara.\\
\(B_m,\ M_m,\ P\) & Tablero residual, máscara de torres y escalado bistocástico de esa máscara.\\
\(\parmar,\ \familias(z)\) & Par de casillas marcado y conjunto de las cuatro líneas de \(z\).\\
\(H_j,\ F_b\) & Llegada histórica corta y llegada sin salida con una completación.\\
\(J,\ S,\ \sumafrontera\) & Contribución de una sola diagonal, magnitud del término del
 cuadrado y suma firmada de frontera \(\sumafrontera=P_T+J-S\).\\
\(\coefT,\ F(x)\) & Primer coeficiente de torres, \(\coefT=\mathfrak F-1/12\),
 y valor de \(\mathfrak F\) en el perfil escalado.\\
\(\centrado,\ \pesodirecto(z)\) & Núcleo de torres centrado y peso de un cruce con sus indicadores
 de disponibilidad.\\
\(d(x),\ \numerr(x)\) & Suma de las dieciséis contribuciones y numerador de la integral del teorema.\\
\(\valorcorte(b)\) & Aproximación al límite con el extremo de integración fijado en \(b\).\\
\bottomrule
\end{tabular}
\refstepcounter{table}\label{tab:notacion}\par\smallskip
\textbf{Tabla \thetable.} Símbolos utilizados en la demostración.
\end{center}
